% !TEX root = ../main.tex
\section{Lower Bound for Constant Center Stepsizes}
\label{sec:lower-bound}
% Source notes (new bibliography entries are deferred to preserve refer.bib):
% Carmon et al., arXiv:1710.11606, Definition 1 and Section 2.3, equation (7).
% Tian--So, arXiv:2210.06907, Section 3.1 and Theorem 3.4.
% Li--Pan, arXiv:2608.03246, Definitions 2.1--2.6 and Remark 3.1.
% Li--Zhu--So, arXiv:2209.10825, Definition 3.1 and Section 7.
% Only the definition formats and the GS/OS distinction are used here.
% In particular, arbitrary new outputs in Li--Pan Remark 3.1 are not covered
% by a hitting-time bound for the triples generated by a fixed recursion.

\subsection{Main result}
\label{subsec:lb-main}

We use the problem class of Definition~\ref{def:problem-class}, the family
\(\calC_\theta\) of Definition~\ref{def:algorithm-family}, and the complexity
of Definition~\ref{def:complexity-measure}, with the normalized constants
\(\mathbf C_0:=(1,1,0,1,1)\). Since \(L_{\calA}=0\), every instance in
\(\mathcal I(\mathbf C_0,\Delta)\) has an affine inner mapping.

\begin{theorem}[Lower bound uniform in \(\rho\) and \(s\)]
\label{thm:lower-bound}
Fix \(0<\eta\le\gamma\), \(\tau>0\), and \(\Delta>0\). There exist constants
\(c_1=c_1(\gamma,\eta,\tau)>0\) and \(c_2=c_2(\gamma,\tau,\Delta)>0\),
independent of \(\rho\), \(s\), \(\alpha\), and \(\beta\), such that the
following holds. For every \(\epsilon\in(0,c_2)\), every \(\rho>0\) and
\(s\ge\tau\) with \(\epsilon(\rho+s)\le c_1\), and every
\((\alpha,\beta)\in[0,1]^2\), there exists an instance
\(I\in\mathcal I(\mathbf C_0,\Delta)\) on which
\(\mathrm{SPLALM}_{\theta,\alpha,\beta}\) requires at least
\[
    \frac{\Delta\,\epsilon^{-4}}{32768\,\gamma\max\{1,4\tau,4\gamma\tau\}}
\]
iterations before one of its outputs is an \(\epsilon\)-KKT point.
Equivalently,
\[T^{\mathrm{KKT}}_\epsilon(\calC_\theta,\mathcal I(\mathbf C_0,\Delta))
\ge\Delta\epsilon^{-4}/(32768\,\gamma\max\{1,4\tau,4\gamma\tau\})\].
\end{theorem}
\todo{State the GS version $T^{\mathrm{GS}}_\epsilon$ after the strict-convex perturbation lemma is proved.}
% perturbation lemma is proved.

\begin{remark}[Initialization]
\label{rem:lb-initialization}
The tight analyses of GDA-type methods
in~\cite[Theorems~5.1--5.2]{mahdavinia2022tight}
and~\cite[Theorems~3.2 and~3.4]{li2026smoothing} choose the hard function
together with the initialization. In Theorem~\ref{thm:lower-bound}, the
initialization is fixed at zero by Definition~\ref{def:splalm}, and only the
instance depends on \((\theta,\alpha,\beta,\epsilon)\).
\end{remark}

\begin{corollary}[Necessary scale of the effective penalty]
\label{cor:effective-penalty}
Let \(\rho=\rho(\epsilon)>0\) and \(s=s(\epsilon)\ge\tau\). If
\(\rho+s=o(\epsilon^{-1})\), then Theorem~\ref{thm:lower-bound} applies for
all sufficiently small \(\epsilon\), and the complexity of \(\calC_\theta\)
is \(\Omega(\epsilon^{-4})\). In particular, an \(O(\epsilon^{-3})\)
guarantee for \(\calC_\theta\) requires \(\rho(\epsilon)+s(\epsilon)>c_1/\epsilon\)
for all sufficiently small \(\epsilon\).
\end{corollary}

\subsection{Relation to Smoothed GDA and DS-GDA}
\label{subsec:gda-embedding}

\(\mathrm{SPLALM}_{\theta,\alpha,\beta}\) in Definition~\ref{def:splalm} uses
\(\prox_{h/\rho}\), so it requires \(\rho>0\).
Smoothed GDA and DS-GDA contain no such term. They nevertheless generate the
same iterates as \(\mathrm{SPLALM}_{\theta,\alpha,\beta}\) as long as
\(\by(\bx,\blmbd)=\zero\). To state this, define
\[
\mathcal R_\rho:=\bigl\{(\bx,\blmbd):\ \blmbd+\rho\calA(\bx)\in\partial h(\zero)\bigr\},
\qquad
\Psi_\rho(\bx,\blmbd):=f(\bx)+\frac{\rho}{2}\|\calA(\bx)\|^2+\langle\blmbd,\calA(\bx)\rangle .
\]

\begin{proposition}[Reduction on \(\mathcal R_\rho\)]
\label{prop:reduction}
Let \(\rho>0\) and \((\bx,\blmbd)\in\mathcal R_\rho\). Then
\(\by(\bx,\blmbd)=\zero\), \(\bq(\bx,\blmbd)=\blmbd+\rho\calA(\bx)\), and
\[
\nabla f(\bx)+J_{\calA}(\bx)^\top\bq(\bx,\blmbd)=\nabla_{\bx}\Psi_\rho(\bx,\blmbd),
\qquad
\calA(\bx)-\by(\bx,\blmbd)=\nabla_{\blmbd}\Psi_\rho(\bx,\blmbd).
\]
Moreover,
\(\calD_\rho(\bx,\blmbd)=\Psi_\rho(\bx,\blmbd)+h(\zero)\).
\end{proposition}

\begin{proof}
Set \(\bv:=\calA(\bx)+\blmbd/\rho\). By~\eqref{eq:prox-multiplier}, \(\by(\bx,\blmbd)\)
is the unique minimizer of \(h(\by)+\frac{\rho}{2}\|\by-\bv\|^2\), and
\(\by\) is this minimizer if and only if \(\rho(\bv-\by)\in\partial h(\by)\).
For \(\by=\zero\) the inclusion reads
\(\blmbd+\rho\calA(\bx)\in\partial h(\zero)\), which holds because
\((\bx,\blmbd)\in\mathcal R_\rho\). Hence \(\by(\bx,\blmbd)=\zero\) and
\(\bq(\bx,\blmbd)=\blmbd+\rho\calA(\bx)\). Since
\(\nabla_{\bx}\bigl(\frac{\rho}{2}\|\calA(\bx)\|^2\bigr)=\rho J_{\calA}(\bx)^\top\calA(\bx)\),
\[
\nabla f(\bx)+J_{\calA}(\bx)^\top\bigl(\blmbd+\rho\calA(\bx)\bigr)
=\nabla_{\bx}\Psi_\rho(\bx,\blmbd),
\]
and \(\calA(\bx)-\zero=\nabla_{\blmbd}\Psi_\rho(\bx,\blmbd)\).
Finally, completing the square in~\eqref{eq:reduced-al} gives
\(\calD_\rho(\bx,\blmbd)=f(\bx)+\min_{\by}\{h(\by)+\frac{\rho}{2}\|\by-\bv\|^2\}
-\frac{1}{2\rho}\|\blmbd\|^2\), and the minimum equals
\(h(\zero)+\frac{\rho}{2}\|\bv\|^2\) because its minimizer is \(\zero\). Hence
\[
\calD_\rho(\bx,\blmbd)
=f(\bx)+h(\zero)+\frac{\rho}{2}\|\calA(\bx)\|^2+\langle\blmbd,\calA(\bx)\rangle
+\frac{1}{2\rho}\|\blmbd\|^2-\frac{1}{2\rho}\|\blmbd\|^2
=\Psi_\rho(\bx,\blmbd)+h(\zero).
\]
\end{proof}

For \(\theta=(\rho,s,\gamma,\eta,\tau)\) and \((\alpha,\beta)\in[0,1]^2\), consider
the recursion
\begin{equation}\label{eq:gda-recursion}
\begin{aligned}
\bx^{k+1}&=\bx^k-\eta\Bigl[\nabla_{\bx}\Psi_\rho(\bx^k,\blmbd^k)
  +\frac{1}{\gamma}(\bx^k-\bz^k)\Bigr],\\
\blmbd^{k+1}&=\blmbd^k+\tau\Bigl[\nabla_{\blmbd}\Psi_\rho(\bx^{k+1},\blmbd^k)
  -\frac{1}{s}(\blmbd^k-\bmu^k)\Bigr],\\
\bz^{k+1}&=\bz^k+\alpha(\bx^{k+1}-\bz^k),\qquad
\bmu^{k+1}=\bmu^k+\beta(\blmbd^{k+1}-\bmu^k).
\end{aligned}
\end{equation}

\begin{proposition}[Embedding]
\label{prop:embedding}
Let \(\rho>0\), and let the sequence
\(\bw^k:=(\bx^k,\blmbd^k,\bz^k,\bmu^k)\) be generated by
\(\mathrm{SPLALM}_{\theta,\alpha,\beta}\). Suppose that for some \(K\ge1\),
\((\bx^k,\blmbd^k)\in\mathcal R_\rho\) and
\((\bx^{k+1},\blmbd^k)\in\mathcal R_\rho\) for all \(0\le k<K\).
Then \eqref{eq:gda-recursion} holds for \(0\le k<K\), and
\[
R_k=\max\bigl\{\|\nabla_{\bx}\Psi_\rho(\bx^{k+1},\blmbd^k)\|,\;
\|\nabla_{\blmbd}\Psi_\rho(\bx^{k+1},\blmbd^k)\|\bigr\},
\qquad 0\le k<K.
\]
\end{proposition}

\begin{proof}
Apply Proposition~\ref{prop:reduction} at \((\bx^k,\blmbd^k)\) in the
\(\bx\)-update of~\eqref{eq:splalm-updates}, and at \((\bx^{k+1},\blmbd^k)\) in
the \(\blmbd\)-update. The same point gives \(\by^{k+1}=\zero\) and
\(\bq^k=\blmbd^k+\rho\calA(\bx^{k+1})\), so the formula for \(R_k\) follows
from~\eqref{eq:kkt-residual}.
\end{proof}

Recursion~\eqref{eq:gda-recursion} is the update of
DS-GDA~\cite{zheng2023universal} for the saddle function \(\Psi_\rho\), with
the same order of updates; the parameter correspondence for DS-GDA and its
special cases is listed in Table~\ref{tab:special-cases}.

\begin{remark}
\label{rem:embedding-scope}
The following points limit the scope of Propositions~\ref{prop:reduction}
and~\ref{prop:embedding}.
\begin{enumerate}[label=(\roman*)]
\item The projections onto \(X\) and \(Y\) in the cited methods are inactive
  when the trajectory stays in the interior of \(X\times Y\). For
  \(h=\|\cdot\|_2\), we have \(\partial h(\zero)=\mathbb B(\zero,1)\) and
  \(h^*=\iota_{\mathbb B(\zero,1)}\), so \(Y=\mathbb B(\zero,1)\) is the
  natural dual set.
\item For a smooth \(\tilde f\), \(\mathrm{SPLALM}_{\theta,\alpha,\beta}\) on the instance
  \(\bigl(\tilde f-\frac{\rho}{2}\|\calA\|^2,\calA,h\bigr)\) generates
  \eqref{eq:gda-recursion} with \(\Psi=\tilde f+\langle\cdot,\calA\rangle\)
  whenever its trajectory stays in \(\mathcal R_\rho\). A hard instance of
  this form for Smoothed GDA or DS-GDA is therefore a hard instance for
  \(\mathrm{SPLALM}_{\theta,\alpha,\beta}\) with any \(\rho>0\), provided
  that \(\tilde f-\frac{\rho}{2}\|\calA\|^2\) satisfies the bounds in
  Definition~\ref{def:problem-class}.
\item Outside \(\mathcal R_\rho\), the two recursions differ. The analysis of
  SPLALM uses the smoothness of \(\calD_\rho(\cdot,\blmbd)\) in~\eqref{eq:reduced-al},
  which requires \(\rho>0\); it does not apply to GDA-type methods.
\item The case \(s=\infty\) is reproduced by \(\beta=1\) for every finite
  \(s\); see Section~\ref{subsec:effective-penalty}. PSGDA with
  \(r_y=\Theta(\epsilon)\) has \(s=\Theta(\epsilon^{-1})\); whether it
  satisfies \(\epsilon(\rho+s)\le c_1\) depends on the constants in \(r_y\)
  and \(c_1\).
\end{enumerate}
\todo{State the transfer of Theorem~\ref{thm:lower-bound} to Smoothed GDA and DS-GDA as a corollary.}
% DS-GDA as a corollary, after the instances in its proof are checked to
% stay in R_rho for all k < T_epsilon.
\end{remark}
